\section{Logarithmic asymptotics of the running inverse coefficient}
\label{sec:running-coefficient}

Let \(G\subset U(r_G)\) be a closed matrix subgroup, in the indicated
unitary representation, with nonabelian simple real Lie algebra
\(\mathfrak g\), where \(r_G\ge1\). Explicitly,
\[
 \mathfrak g=\{X\in M_{r_G}(\mathbb C):
       \exp(tX)\in G\text{ for every }t\in\mathbb R\}.
\]
Since \(U(r_G)\) is closed in \(\mathrm{GL}(r_G,\mathbb C)\), so is \(G\).
Thus Hall's matrix Lie group results apply: this set is a finite-dimensional
real vector space closed under \([X,Y]=XY-YX\)
\cite[Definition~3.13 and Theorem~3.16]{hall2000}.
Since \(\exp(tX)\) is unitary for
\(X\in\mathfrak g\) and real \(t\), differentiating
\(\exp(tX)^*\exp(tX)=I\) at \(t=0\) gives \(X^*=-X\).
Thus the symmetric real bilinear form
\(\langle X,Y\rangle=-r_G^{-1}\operatorname{Re}\operatorname{tr}(XY)\)
is positive definite: for \(X\ne0\),
\[
 \langle X,X\rangle
 =r_G^{-1}\operatorname{tr}(X^*X)
 =r_G^{-1}\sum_{a,b}|X_{ab}|^2>0.
\]
Nonabelianness implies \(\mathfrak g\ne\{0\}\), so
\(\dim_{\mathbb R}\mathfrak g\ge1\). Write
\[
    b_0(G)=\frac{11\,C_{\mathfrak g}}{48\pi^2},
    \qquad
    C_{\mathfrak g}
      =-\frac{1}{\dim_{\mathbb R}\mathfrak g}
        \sum_j \operatorname{Tr}_{\mathfrak g}
             \bigl((\operatorname{ad}e_j)^2\bigr),
\]
where \((e_j)\) is an orthonormal basis for this inner product and
\(\operatorname{ad}X(Y)=[X,Y]=XY-YX\). Here
\(\operatorname{Tr}_{\mathfrak g}\) is the real trace of an endomorphism of
\(\mathfrak g\). This constant is independent of the orthonormal basis:
if \(f_i=\sum_j q_{ij}e_j\) with \(q\) orthogonal, then
\[
 \sum_i(\operatorname{ad}f_i)^2
 =\sum_{j,k}\Bigl(\sum_iq_{ij}q_{ik}\Bigr)
       \operatorname{ad}e_j\operatorname{ad}e_k
 =\sum_j(\operatorname{ad}e_j)^2.
\]
Taking traces proves the assertion. Invariance of the matrix trace gives
\(\langle[X,Y],Z\rangle=-\langle Y,[X,Z]\rangle\), so each
\(\operatorname{ad}e_j\) is skew-adjoint. Consequently
\(-\operatorname{Tr}_{\mathfrak g}((\operatorname{ad}e_j)^2)
=\sum_\ell\|[e_j,e_\ell]\|^2\). If every basis commutator vanished,
bilinearity would make all commutators vanish. Noncommutativity therefore
makes at least one term in the resulting double sum positive. In particular,
\(b_0(G)>0\).

\paragraph{Input from the lattice construction.}
Fix a construction family \(\mathcal C\) satisfying the construction
conditions. For \(N,p,m\in\mathbb N\) with \(m\le N\), let
\(\beta_{N,p,m}\) be its action coefficient at stage \(m\) of the
trajectory started at lattice level \(N\) and period index \(p\).
The spacing at stage \(m\) is \(a_m=(\prod_{j<m}B_j)^{-1}\), where each
block factor \(B_j\) is an integer at least \(3\). Thus
\(s_m=\log(a_m^{-1})=\sum_{j<m}\log B_j\ge m\log 3\), so \(s_m\to+\infty\).
The running estimate for the family provides a fixed reference intercept
\(\bar\beta\) and, for every
\(\eta>0\), a threshold \(m_0(\eta)\) such that
\begin{equation}
 \left|\beta_{N,p,m}-\bigl(\bar\beta+2b_0(G)s_m\bigr)\right|
 \leq \eta s_m
 \qquad
 \bigl(m\geq m_0(\eta),\ N\geq m,\ p\in\mathbb N\bigr).
 \label{eq:running}
\end{equation}
The threshold is uniform over the starting level and period index.
We use the running estimate as an input to the comparison below.

\begin{proposition}[Uniform extraction]
\label{prop:uniform-extraction}
Let $s_m\to+\infty$, and let $A,c\in\mathbb R$ be fixed. Suppose the real
coefficients $\beta_{N,p,m}$ satisfy, for every $\eta>0$, an estimate
$|\beta_{N,p,m}-(A+cs_m)|\le\eta s_m$ for all sufficiently large $m$,
uniformly over $N\ge m$ and $p\in\mathbb N$.
Then for every $\varepsilon>0$, for all sufficiently large $m$ one has $s_m>0$
and
\[
 \left|\frac{\beta_{N,p,m}}{s_m}-c\right|<\varepsilon
 \qquad(N\ge m,\ p\in\mathbb N).
\]
\end{proposition}

\begin{proof}
Fix $\varepsilon>0$. Use the assumed estimate with $\eta=\varepsilon/2$,
and enlarge its threshold so that $s_m>2(|A|+1)/\varepsilon>0$.
Then, for every $N\ge m$ and $p$,
\[
 \left|\frac{\beta_{N,p,m}}{s_m}-c\right|
 \le \frac{|\beta_{N,p,m}-(A+cs_m)|}{s_m}+\frac{|A|}{s_m}
 \le\frac\varepsilon2+\frac{|A|}{s_m}<\varepsilon.
\]
Both thresholds are independent of $N,p$. No positivity assumption on $c$
is used in this proposition.
\end{proof}

\begin{proposition}[Logarithmic asymptotic]
\label{prop:running-asymptotic}
Let \(N,k:\mathbb N\to\mathbb N\) be arbitrary sequences with \(i\leq N(i)\)
for all sufficiently large \(i\). For \(\beta_{N,p,m}>0\) set
\(g_{N,p,m}=(\sqrt{\beta_{N,p,m}})^{-1}\). Then
\(\beta_{N(i),k(i),i}>0\) for all sufficiently large \(i\), and
\[
  \lim_{i\to\infty}g_{N(i),k(i),i}\sqrt{\log(a_i^{-1})}
  =\frac{1}{\sqrt{2b_0(G)}}.
\]
\end{proposition}

\begin{proof}
Apply Proposition~\ref{prop:uniform-extraction} to \eqref{eq:running}, with
$A=\bar\beta$ and $c=2b_0(G)>0$. For each $\varepsilon>0$, enlarge the
resulting threshold until $i\le N(i)$ holds. Substituting
$m=i$, $N=N(i)$ and $p=k(i)$ gives
\[
 \left|\frac{\beta_{N(i),k(i),i}}{s_i}-c\right|<\varepsilon
\]
for every subsequent $i$. Thus $x_i=\beta_{N(i),k(i),i}/s_i\to c$ on the
tail where $s_i>0$. Taking $\varepsilon=c/2$ shows $x_i>c/2>0$ eventually;
hence $\beta_{N(i),k(i),i}=s_ix_i>0$ there. No condition on the sequence $k$
was needed, because the original threshold was uniform in $p$.

On this positive tail the square-root limit rule
\cite[Proposition~2.2.6]{lebl} gives $\sqrt{x_i}\to\sqrt c$.
All denominator terms and their limit are nonzero, so the quotient rule
\cite[Proposition~2.2.5(iv)]{lebl}, with constant numerator $1$, gives
$1/\sqrt{x_i}\to1/\sqrt c$. Finally,
\[
 \frac1{\sqrt{x_i}}
 =\frac{\sqrt{s_i}}{\sqrt{\beta_{N(i),k(i),i}}}
 =g_{N(i),k(i),i}\sqrt{\log(a_i^{-1})}.
\]
The first identity follows by squaring the positive quantities
$\sqrt{x_i}$ and $\sqrt{\beta_{N(i),k(i),i}}/\sqrt{s_i}$ and using uniqueness
of the nonnegative square root. All expressions are interpreted on the
positive tail and may be assigned arbitrary values at earlier indices:
changing finitely many terms preserves convergence by enlarging each tolerance's
threshold beyond those indices. This proves the stated limit.
\end{proof}

At \(N(i)=i\) the coefficient is the bare lattice action coefficient.
For a general admissible \(N(i)\), the comparison concerns the coefficient
read from the trajectory at spacing \(a_i\). Relating this inverse-square-root
proxy to a coupling defined through a physical observable would require a
separate matching theorem.
